\labsection{5}{Design graphics that answer a question}{sec:lab05}

\begin{labproject}{Visual communication: turn analysis into evidence}
\textbf{Coverage.} Chapter~\ref{ch:visualization}.\\
\textbf{Goal.} You should be able to map analytical questions to chart forms, identify misleading encodings, write evidence-focused captions, and distinguish information revealed by raw and aggregated views.\\
\textbf{Datasets.} \texttt{red\_wine\_quality\_dataset.csv}; Bike Sharing hourly data are used for the extension.\\
\textbf{Scenario.} You are preparing figures for an analytical report. The objective is not decoration: each chart must answer a specific question, use an appropriate visual encoding, and make the intended comparison easy to see.

\textbf{Core design tasks.}
\begin{enumerate}
  \item \textbf{Distribution question.} Use a bar chart to show the frequency of the discrete wine-quality scores. Explain why a histogram is not the clearest default for this variable.
  \item \textbf{Group-comparison question.} Use side-by-side boxplots to compare alcohol content across quality scores. Support the figure with grouped numerical summaries.
  \item \textbf{Relationship question.} Use a scatterplot to examine volatile acidity and alcohol. Describe what can and cannot be inferred when neither variable is an outcome in the graphic.
  \item \textbf{Redesign by changing the question.} Create a second version of one figure that answers a genuinely different analytical question. Do not count changes in colour, font, or theme as a redesign.
  \item \textbf{Write analytical captions.} For each figure, state the data subset, visual encoding, main visible pattern, and one limitation or uncertainty.
  \item \textbf{Perform a visual integrity check.} Identify one choice that would exaggerate, hide, or confuse the pattern---for example, an inappropriate axis, unnecessary 3-D effect, excessive overplotting, or treating a discrete score as continuous without explanation.
\end{enumerate}

\begin{tryit}
\textbf{Extension---two-dimensional time structure.} Use \texttt{bike\_sharing\_dataset\_hour.csv} to create an hour-by-weekday heatmap of mean rentals. Explain what the heatmap reveals about commuting structure that a single hourly line cannot, and identify what variability is hidden by averaging.
\end{tryit}

\tcbbreak
\lstinputlisting[style=utpr]{all_experiments/lab05_visualization.R}
\begin{verification}
\textbf{Required evidence.} Submit a three-figure mini-report. Each figure must have a one-sentence question above it and a two-to-three sentence evidence caption below it. A reader should understand the analytical claim without seeing the R code.

\textbf{Professional standard.} Prefer the simplest chart that makes the required comparison accurately. Visual complexity is justified only when it reveals additional structure.
\end{verification}
\end{labproject}
